\documentclass[12pt,letterpaper]{article}
\usepackage[T1]{fontenc}
\IfFileExists{newtxtext.sty}{\usepackage{newtxtext,newtxmath}}{\usepackage{mathptmx}\usepackage{amssymb}}
\usepackage[letterpaper,margin=1in]{geometry}
\usepackage{amsmath,booktabs}
\usepackage{microtype}
\usepackage{enumitem}
\usepackage[hidelinks]{hyperref}
\title{Revision log: genuine invariant distributions on metaplectic covers of $SL(2)$}
\author{David Joyner\\\texttt{wdjoyner@gmail.com}}
\date{Drafts 2 and 3 --- November 21, 2026}
\begin{document}
\maketitle
\section*{Sources, scope, and decision}
The inputs were the attached January 4, 1994 manuscript
\texttt{hitchhiker-sl2bar-repns.tex}, the attached corrected 2026d source
\texttt{joyner\_invar-distrib\_corrected-draft2026d.tex}, and the attached
mathematical-paper style guide (November 11, 2026).  The revised manuscript
adopts the original 1994 title, with a new 2026 date line.  It is a
\emph{synthesis}, not a line-for-line edited republication of either the
1994 manuscript or the 2001 journal article.  Material on $GL(r)$ is omitted
except for bibliographic provenance.

\section*{Changes from the supplied corrected 2026d draft}
\begin{enumerate}[leftmargin=*,itemsep=4pt]
\item Replaced the 2001 article's title with the exact title of the 1994
metaplectic manuscript.  Identified both antecedents in a title footnote.
Recorded 1994 as January 4 (the original has \texttt{1-4-94}) and the
published 2001 article separately.
\item Restricted the mathematical setting to $SL(2,F)$, with $F$ a finite
extension of $\mathbb Q_p$, $\mu_n\subset F$, and $p\nmid 2n$.
This is narrower than the 1994 manuscript and is not a result for wild covers.
\item Rewrote the abstract and introduction; added an explicit split
between algebraic trace density and the continuity/topology problem.
\item Inserted literature reviews for the sources through 1994 and
post-1994 developments, with a direct discussion of Li's 2012 Sections~5.4,
5.6, and Theorem~5.8.10.
\item Added the central-extension group law from the 1994 paper, and
several elementary examples clarifying genuine test functions, tame fields,
conjugation differences, and spectral distributions.
\item Reorganized the notation table into a stand-alone early section.
Removed scope-specific assumptions, comments and claims about covers of
$GL(r)$.  Preserved the older internal theorem labels when possible.
\item Replaced the cited maximal-abelian subgroup lemma by a self-contained
proof of the index-$\nu$ subgroup statement and separated the stronger
maximal-abelian assertion as source-dependent.
\item Identified the orbital-integral surjectivity statements and proposed
principal-series Fourier-image formulas as source-dependent and in need of
verification.  This prevents an unverified surjectivity assertion from
being mistaken for a new proof.
\item Rewrote the kernel-of-trace proof as a deduction from the density
result in Li, Theorem~5.8.10, subject to the hypotheses and representation
inputs in that source; did not claim to reprove Li's theorem.
\item Identified a material conflict: the 1994 typescript calls the image of $J_\mu$
the special constituent and the kernel non-tempered, whereas the supplied
2026d correction reverses them.  Demoted the endpoint proposition to a
verification remark and inserted a separate section surveying the original
conjectures without claiming they were all resolved.
\item Clarified that the quotient-dual isomorphism is algebraic;
continuity is an additional matter.
\item Rewrote the final integral theorem as a \emph{conditional}
continuous Paley--Wiener spectral pairing, with explicit kernel and
quotient-topology requirements.  Removed the assertion that an arbitrary
spectral functional is given by an ordinary density $m(\pi)\,d\mu(\pi)$.
\item Clarified the finite-order definition, retained the Fourier-series
proof of the Schwartz lemma, and restricted its equivariant clause to an
isometric group action.
\item Removed the in-body revision commentary. Added an end-of-paper
summary, open verification items, and glossary.  Updated acknowledgments,
removed the obsolete USNA address and retained the user's current email.
\item Applied 12pt, US Letter, one-inch margins, Times text and math,
and \texttt{pdflatex} compatibility.  Converted bibliography keys to
author-initial/year labels where documentary years were available.
\end{enumerate}

\section*{Draft 2 editorial corrections (November 21, 2026)}
\begin{enumerate}[leftmargin=*,itemsep=4pt]
\item Replaced the temporary ``Draft 2026e'' designation with \emph{Draft 2} throughout the new manuscript and this log. Under the trace-formulas convention, filenames do not carry draft numbers; Git history records successive drafts.
\item Corrected the first sentence of the abstract to disclose AI assistance by name (ChatGPT and Claude), identify the 1994 manuscript, and refer to the Acknowledgments. Removed the redundant AI-disclosure sentence at the end of the abstract.
\item Established the intended repository location as \texttt{sl2bar-invariant-distributions/} inside \texttt{wdjoyner/trace-formulas}.
\item No mathematical statements or proofs were changed during this editorial correction.
\end{enumerate}


\section*{Draft 3 (November 21, 2026): response to two referee reports}
\paragraph{Inputs and AI usage.}
Draft~2 was refereed by Gemini (Google), which the author calls Carbon (file
\texttt{referee-2026-11/report-gemini.md}), and by Claude (Anthropic), who also
assessed Gemini's report item by item (\texttt{referee-2026-11/report-claude.md}).
Claude then prepared Draft~3 in response to both reports.  Claude wrote the new
proofs listed below and the script \texttt{scripts/genuine\_torus\_check.py}.
Draft~3 has not been refereed.

\paragraph{Correction to Claude's own report.}
Running the new script showed that item~9 of Claude's report was wrong in one
respect.  The cocycle $\beta$ is trivial on $N\times N$ for every unipotent radical
$N$, including the lower triangular one, by the Steinberg relation.  So Draft~2's
identification of $N$ with $\{(u,1)\}$ was correct.  Draft~3 keeps it and adds the
one-line reason.  Claude's report also proposed moving the notation table to an
appendix.  That conflicts with the style guide (\S2, item~5, notation near the
beginning), so the table was kept in place.

\paragraph{Results newly established in Draft 3.}
These are new mathematics, not present in the 1994 notes, the 2001 article, or Draft~2.
\begin{enumerate}[leftmargin=*,itemsep=2pt]
\item Lemma~5.2 (formerly Remark~5.2, source-dependent): $\beta(h(a),h(b))=(b,a)_n$, the
commutator in $\bar A$ is $(a,b)_n^{-2}$, the center of $\bar A$ is the preimage of
$A^\nu$ of index $\nu^2$, and $\overline{h(C)}$ is maximal abelian of index $\nu$.  Proved,
and checked exactly by the script for six pairs $(p,n)$ with $p\le19$.
\item Lemma~5.3 (formerly Remark~5.3(a)): vanishing of the split orbital integral off the
center of $\bar A$.  Proved.
\item Lemma~6.3 (restated): the duality pairing is now sesquilinear, between $I(\mu)$ and
$I(\bar\mu^{-1})$.  The Draft~2 bilinear pairing between $I(\mu)$ and $I(\mu^w)$ is not
well defined for $n>2$, because the contragredient of a genuine representation is not
genuine.  Well-definedness and invariance are proved; non-degeneracy is cited.
\item Necessity of the Weyl-fixed hypothesis in Proposition~6.6 (proved in the text after it).
\item Proposition~6.8 (formerly Remark~6.8, an unresolved conflict): for $\mathrm{Re}\,s>0$,
$\ker T_{w_0}$ is the unique irreducible subrepresentation and $\mathrm{im}\,T_{w_0}$ the
Langlands quotient.  At the endpoint, $\ker T_{w_0}\simeq\pi_{sp}$ and
$\mathrm{im}\,T_{w_0}\simeq\pi_{nt}$.  This agrees with the corrected 2026 draft and
reverses the 1994 conjecture.  The proof uses the geometric lemma, Frobenius reciprocity,
and the Langlands quotient theorem for covers (Ban--Jantzen 2013).
\end{enumerate}

\paragraph{Corrections of errors in Draft 2.}
\begin{enumerate}[leftmargin=*,itemsep=2pt]
\item Harish-Chandra's $\Xi$ (\S3.2): the integral of $\delta_B(xk)^{-1/2}$ collapses to
$\delta_B(x)^{-1/2}$.  Replaced by $\int_{K_0}\delta_B(kx)^{1/2}\,dk$.  (Claude's report.)
\item Unramified characters of $\bar A$ and the parameter space: $A^n$ replaced by $A^\nu$, the
center (\S3.2, \S8.1).  (Gemini.)
\item Characters $\Theta_\pi$ are anti-genuine, not genuine, under the paper's convention (\S3.2).  (Claude.)
\item $w_0$ has order four (\S3.2).  (Claude.)
\item Splitting over $K_0$: citation changed from Gelbart (the case $n=2$) to Kazhdan--Patterson,
with the note that the splitting is not $k\mapsto(k,1)$.  (Claude.)
\item \S6.2: genuine characters of $\bar C$ written as $\mu_0\eta_s$, with $\eta_s$ pulled back from $A$,
in place of ``$\mu(x)=\mu_0(x)|x|^s$''.  Defined $I(\mu):=I(\chi_\mu)$.  Used $\delta_B^{1/2}$ in $V(\mu)$.
$K(\mu)$ is now defined as the space of restrictions to the big cell, and condition (ii) is
replaced by the exact formula from the Bruhat identity.  The decay is at $|x|\to\infty$ (noncompact
picture); Gemini's suggestion of $|x|$ small was not adopted.  Named $T=T_{w_0}$.  (Gemini and Claude.)
\item Lemma~6.4 is stated for the Hermitian form $H=\langle\cdot,\iota T\cdot\rangle$, which fixes the
bilinear/sesquilinear mismatch.  (Gemini; root cause found by Claude.)
\item Lemma~6.5: range $0<\mathrm{Re}\,s$ (the integral does not converge at $\mathrm{Re}\,s=0$), and
the statement is about the image of $T_{w_0}$, not of the matrix $J_\mu$.  (Claude.)
\item Proposition~6.6 and Corollary~6.7: added the hypotheses $s$ real, $\chi_{\mu_0}$ Weyl-fixed, and
$I(\mu_0)$ irreducible.  (Claude.)
\item Theorems~6.9 and 6.10: the irreducible $I(\chi)$ with $\chi^w=\chi$ and $n$ even was missing.
``The special representation'' changed to ``a special representation''.  (Claude; Gemini.)
\item $\Pi(\bar C)_u$ is a countable disjoint union of circles, not a compact torus (\S8.1).  (Claude.)
\item Theorem~8.6: $\mathcal F(C_c^\infty)=PW_t$ by definition, so ``$\subset$'' became ``$=$''.  (Claude.)
\item Example~8.12 claimed that the conditional Theorem~8.8 implies the Plancherel statement.  It
was merged into Example~9.1 with conditional wording.  (Claude.)
\end{enumerate}

\paragraph{Changes of status.}
Remark~5.2 $\to$ Lemma~5.2 (proved).  Remark~5.3 $\to$ Lemma~5.3 (part (a), proved) and Remark~5.4
(parts (b), (c), still source-dependent).  Remark~6.8 $\to$ Proposition~6.8 (proved, using
(Ban--Jantzen 2013)).  Theorems~6.9 and 6.10 $\to$ Remarks~6.9 and 6.10 (source-dependent classification;
Draft~2 already said they were not verified).  No status was upgraded without a proof.

\paragraph{Exposition and housekeeping.}
Title page and date line in the form of the style guide (\S6.2).  Abstract names all three AI systems.
Introduction states the tame restriction $p\nmid2n$ and mentions Section~7.  Literature review adds
Ban--Jantzen (2013, 2016), Kaplan--Szpruch (arXiv 2022) and the 2024 corrections to Li's paper.
\S8 opening softened from ``we classify'' to ``we review''.  Note that algebraic and continuous duals
of $C_c^\infty$ coincide.  The duplicate definition of the space of conjugation differences removed.  \S8.1 rewritten in modern
\LaTeX.  The \verb|\F| alias and the \verb|\verify|, \verb|\cD|, \verb|\aaa| macros and the unused
cleveref package removed.  Commented verification notes turned into text.  Proofs use the
\texttt{proof} environment.  ``2026d'' replaced by a citation ([Jo26] in the paper).  Bibliography typos fixed
(centre, corps, Zelevinsky, non-Archimedean, forms, Th\'eorie, int\'egrales).  Uncited items Ar81, Be75,
Da99, HC84, Jo88, JM00 and Me98 removed.  Acknowledgments name each AI and its role and point to this log.
Notation table updated.  The preamble falls back to mathptmx when newtx is not installed (the
committed PDF was built that way, because newtx was not available in Claude's workspace).

\paragraph{Verification levels (style guide \S7.3).}
Level~1 (independent derivation): Lemmas~5.2, 5.3, 6.3 (except non-degeneracy); the necessity
argument after Proposition~6.6; Proposition~6.8 given its cited inputs; the Bruhat identities and the
unipotent computation (also checked by the script).  Level~2 (source consulted): only the abstract and
introduction of Ban--Jantzen (2013) and the title and authors of Kaplan--Szpruch.  Level~3 (memory,
not consulted): Arit\"urk, Flicker--Kazhdan, Moen, Kazhdan--Patterson (splitting over $K_0$),
Ban--Jantzen (2016), Vign\'eras, and the content of the 2024 corrections to Li's paper.

\paragraph{Sources not consulted.}
None of Arit\"urk (1980), Flicker--Kazhdan (1986), Moen (1988), or the body of Ban--Jantzen was read
in preparing Draft~3.  The normalization of the endpoint $s=1/n$ and Moen's exceptional datum remain
unchecked.

\paragraph{Edits of November 22, 2026, at the author's request.}
Statements that were present only because of errors in earlier drafts were removed or shortened.
Removed: the remark on the wrong form of $\Xi$ (\S3.2); the sentence on the order of $w_0$ (\S3.2); the
remark contrasting $K(\mu)$ with the Kirillov model (\S6.2); the sentence that no intrinsic description of
$PW_t$ is proved (\S3.2); the remark that a single partial derivative does not suffice in the Schwartz
lemma (\S8.3).  Shortened: the justification that $u\mapsto(u,1)$ is a homomorphism on unipotent radicals
(\S3.2) and the note on Moen's exceptional datum (\S6.1), to one or two sentences each; the paragraph after
Lemma~6.3 on the bilinear pairing, to one sentence; the note that the splitting over $K_0$ is not
$k\mapsto(k,1)$, to a parenthetical remark.  Two disclaimers about scalar densities (the opening of \S8 and
Theorem~8.6) became footnotes.  No mathematical statement changed.  Theorem numbers are unchanged.

\paragraph{Concordance, Draft 2 $\to$ Draft 3.}
\begin{center}\small
\begin{tabular}{@{}ll@{}}
\toprule
Draft 2 & Draft 3\\
\midrule
Remark 5.2 (maximal abelian preimages) & Lemma 5.2 (commutators in the covering torus)\\
Remark 5.3 (orbital integrals) & Lemma 5.3 (part (a)) and Remark 5.4 (parts (b), (c))\\
Example 5.4 (discriminant) & Example 5.5\\
Lemma 6.5 (unlabeled) & Lemma 6.5 (label \texttt{lem:image-irred})\\
Corollary 6.7 (unlabeled) & Corollary 6.7 (label \texttt{cor:positive})\\
Remark 6.8 (endpoint conflict) & Proposition 6.8 (endpoint constituents)\\
Theorem 6.9, Theorem 6.10 & Remark 6.9, Remark 6.10\\
Lemma 8.5 (unlabeled) & Lemma 8.5 (label \texttt{lem:quotient-dual})\\
Example 8.12 & merged into Example 9.1\\
\bottomrule
\end{tabular}
\end{center}
All other numbers are unchanged, and all Draft~2 labels are kept.

\section*{Mathematical verification status}
\begin{description}[style=nextline]
\item[Directly checked] Algebraic quotient argument for invariant
functionals; convolution invariance of traces; finite-order torus lemma
via Fourier series; distinction between ordinary densities and derivative
distributions.
\item[Checked in a cited source] W.-W. Li (2012), Theorem~5.8.10:
vanishing of genuine tempered characters is equivalent to vanishing of
regular orbital integrals and membership in the span of conjugation
differences.  Li, Theorem~5.4.1, describes an $R$-group decomposition
admitted as a hypothesis in the non-Archimedean case.  The theorem is
used as a cited input, not independently reproved.
\item[Conditional] Existence and precise topology of the full genuine
tempered Paley--Wiener space underlying the distributional pairing.
\item[Proved in Draft 3] Lemmas 5.2, 5.3 and 6.3 and Proposition 6.8; see the Draft~3 section above.
\item[Unverified against original sources] The normalization $s=1/n$ of the endpoint
in the dual classification; Moen's exceptional genuine character;
intertwining-form normalizations; the proposed orbital-integral image
and principal-series spectral-image surjectivity; some bibliographic
items originally supplied as unpublished preprints.
\end{description}
\section*{Relation to earlier numbering}
The source theorem labels \texttt{thm:trace}, \texttt{thm:main},
\texttt{prop:kernel}, \texttt{lemma:schwartz}, and the primary section
labels were retained, although their printed numbers shifted after
inserting the new introduction and literature/background material.
An exact 1994-to-2026 numbered theorem concordance is not yet available:
the 1994 typescript is 2,282 lines and this editorial synthesis has not
established statement-by-statement identity.  Do not treat coincident
numbers as equivalent results.

\section*{Publication and copyright caution}
The 2001 journal article is published material and may be subject to
publisher-specific reuse conditions. The present standalone expository
synthesis is not a reproduction or replacement of the published article.
The author requested that Draft 2 be posted in the public
\texttt{wdjoyner/trace-formulas} repository, in
\texttt{sl2bar-invariant-distributions/}.
\end{document}
